\documentclass[11pt,a4paper]{article}
\usepackage[margin=1in]{geometry}
\usepackage[T1]{fontenc}
\usepackage{lmodern}
\usepackage{amsmath,amssymb}
\usepackage{booktabs,tabularx,array}
\usepackage{caption}
\usepackage{microtype}
\usepackage[hidelinks]{hyperref}
\usepackage{enumitem}
\hypersetup{
  pdftitle={Evaluating Direct and Solver-Assisted AI Degree Planning},
  pdfauthor={Lu Ruijian; Brandon Chong Wee Kiat; Jonathan Graziano Yang}
}
\captionsetup{font=normalsize,labelfont=bf}
\setlength{\parindent}{0pt}
\setlength{\parskip}{0.25em}
\setlength{\emergencystretch}{2em}
\setlist{nosep,leftmargin=*}
\newcolumntype{Y}{>{\raggedright\arraybackslash}X}
\newcommand{\code}[1]{\texttt{#1}}
\newcommand{\course}[1]{\texttt{#1}}
\renewenvironment{abstract}
  {\section*{Abstract}\begingroup\normalsize}
  {\endgroup}

\begin{document}
\hypersetup{pageanchor=false}
\begin{titlepage}
\centering
\vspace*{1.2cm}
{\Large National University of Singapore\\School of Computing\par}
\vspace{1.5cm}
{\LARGE\bfseries Evaluating Direct and Solver-Assisted\\[0.25em]
AI Degree Planning\par}
\vspace{0.7cm}
{\Large Schedule Validity, Translation Fidelity\\
and Explanation Accuracy\par}
\vspace{1.2cm}
{\large CS3263 Foundations of Artificial Intelligence\\
Semester 1, AY2026/27\par}
\vspace{1.3cm}
\begin{tabular}{lll}
\toprule
Team member & Matriculation number & Email\\
\midrule
Lu Ruijian & A0287930U & e1297772@u.nus.edu\\
Brandon Chong Wee Kiat & A0269806N & e1114718@u.nus.edu\\
Jonathan Graziano Yang & A0300413A & e1354779@u.nus.edu\\
\bottomrule
\end{tabular}
\vfill
{\large Final project report\\11 November 2026\par}
\end{titlepage}
\pagenumbering{arabic}
\hypersetup{pageanchor=true}

\begin{abstract}
We compare direct language-model planning with translation into a constraint
satisfaction problem (CSP) and symbolic solving. We study remaining NUS
Computer Science Programme Requirements with the AI focus, a restricted
catalogue and a 16-unit semester cap. The frozen evaluation contains 120
cases: 84 feasible and 36 infeasible across four horizons and three designed
difficulty levels. Direct planning achieved 119/120 correct
outcomes; solver assistance achieved 120/120. Mean recorded process plus
solver time was 22.41 and 31.60 seconds respectively. All five noncanonical
translations were reviewed as semantically equivalent. A secondary audit
identified 17 incorrect numeric claims in 15 otherwise valid direct-plan
responses. Both methods achieved near-ceiling outcome accuracy on these
structured instances. Separating schedule validity, translation fidelity
and explanation accuracy exposed errors concealed by the primary score.
\end{abstract}

\section{Introduction and research question}
Degree planning coordinates course selection, prerequisite order, recurring
offerings and workload. A course combination may meet aggregate requirements
while remaining impossible to schedule; a plausible schedule may omit
requirements or miscount overlapping groups. These interactions motivate a
constraint satisfaction model for an educational assistant.

The project asks: \emph{under the tested conditions, how reliably does a frontier
language model plan directly, and what changes when it translates the request
for a formal solver?} We examine three proposal sub-problems: selecting courses
that satisfy requirements, respecting prerequisite relationships, and assigning
courses to semesters subject to workload limits. Soft preferences and
optimization are deferred. Every solution that satisfies the specified hard
constraints is acceptable.

We contribute an explicit Boolean CSP, a benchmark with validated outcomes,
and a paired evaluation separating schedule validity, translation fidelity
and explanation accuracy. Prior work couples language models with classical
planning in LLM+P \cite{llmp} and with symbolic inference in Logic-LM
\cite{logiclm}. We study this division of labour in a restricted degree-planning
domain, with a fixed compiler and one attempt per request. Our contribution
is the checked benchmark and case-level comparison, including inaccuracies
in explanations of valid schedules.

\section{Domain and scope}
The domain is the AY2026/27 NUS BComp(Computer Science) admission cohort and AI
focus area. Official programme and focus-area pages provide the requirements;
NUSMods supplies course records \cite{nuscohort,nusfocus,nusmods}. Source
snapshots were collected on 9 October 2026. The normalized catalogue contains
87 courses: nine foundations, three mathematics courses, 19 AI breadth
courses, one credited internship and 55 prerequisite-support courses.
CP3200 is completed-credit-only. CS2040HS, EE2028 and EE2211 are conservatively
excluded from new enrolment because wildcard prerequisites remain unresolved
by the normalizer; these exclusions reflect parsing limitations.

The nominal 80-unit Programme target combines 36 Foundation units,
12 Mathematics units and at least 32 Breadth/Depth units. The model excludes Common
Curriculum, Unrestricted Electives, GPA, exemptions, exchange mappings,
dissertation substitutions, special terms, quotas and timetable clashes.
Prerequisite-support courses consume workload without earning Programme
credits. Benchmark requests explicitly require CS1101S although its credits
are outside Programme groups.

The workload limit of 16 modelled units per normal semester is a project
policy, reserving capacity for other study; it is not an official NUS maximum.
Empty semesters are permitted. Normal Semesters 1 and 2 alternate, and the
AY2026/27 offering pattern repeats in simulated future years. Standard
mathematics, English and programme-entry conditions are assumed satisfied;
passed courses satisfy modelled grade thresholds. CP3200 contributes six
completed internship units in every case, including early-stage profiles.
Internship eligibility is not simulated. These choices define a simplified
remaining-requirements planner.

\subsection{A representative planning request}
In evaluation case \code{test-h4-easy-01}, the completed record includes
CS1101S and the student has four remaining semesters, starting with normal
Semester 1. The request contains the exact excerpt:
\textquotedblleft My plan must include CS1101S, CS4277.\textquotedblright\
Inclusion can be satisfied by past completion; only CS4277 needs new
enrolment among these two requested courses. Table~\ref{tab:worked} shows the
saved assisted schedule, which passes the original-specification checker.

\begin{center}
\begin{minipage}{\linewidth}
\centering
\captionof{table}{A checked schedule for the representative request.}
\label{tab:worked}
\begin{tabularx}{\linewidth}{@{}rrYr@{}}
\toprule
Position & Normal semester & New enrolments & Units\\
\midrule
1 & 1 & CS2100, CS3268, CS4261 & 12\\
2 & 2 & CS2106, CS4243, CS4248 & 12\\
3 & 1 & None & 0\\
4 & 2 & CS3263, CS4277, CS5242 & 12\\
\bottomrule
\end{tabularx}
\end{minipage}
\end{center}
CS2100 precedes CS2106, CS4277 is offered in the assigned normal semester,
and CS1101S is not repeated. With no minimum load, semester 3 may remain
empty; no optimization objective is applied.

\par
\section{Formal CSP formulation}
\subsection{Inputs, variables and completed credit}
Let $C$ be the frozen catalogue, $T=\{1,\ldots,H\}$ the remaining semester
positions, and $a_t\in\{1,2\}$ the normal-semester number at position $t$.
For course $c$, let $u_c$ denote its units, $A_c$ its offering set and
$q_c\in\{0,1\}$ its new-enrolment eligibility. Let $E\subseteq C$ contain
completed courses, with $d_c=\mathbf{1}[c\in E]$. Request inputs specify
required courses $R$, forbidden new enrolments $F$, fixed placements
$J\subseteq C\times T$, workload cap $W$, and requirement groups.

The CSP is $\mathcal{P}=(V,D,K)$, with finite Boolean domains and hard
constraints $K$. Its principal variables are
\[
 x_{c,t}\in\{0,1\},\qquad y_c\in\{0,1\}.
\]
Here $x_{c,t}=1$ denotes new enrolment at remaining position $t$; $y_c=1$
denotes a course completed previously or selected in the plan. Enforce
\begin{align}
 \sum_{t\in T}x_{c,t}&\leq 1, \label{eq:once}\\
 \sum_{t\in T}x_{c,t}&=0
       &&\text{if }d_c=1\text{ or }q_c=0, \label{eq:elig}\\
 y_c&=d_c+\sum_{t\in T}x_{c,t}. \label{eq:link}
\end{align}
Thus a completed required course satisfies $y_c=1$ without being enrolled
again. A non-selectable course may still contribute trusted completed credit.
Forbidden courses retain past credit but cannot be newly taken.

\subsection{Offerings, prerequisites and co-requisites}
Offering restrictions impose $x_{c,t}=0$ whenever $a_t\notin A_c$.
For a prerequisite course $p$, define completion before $t$ by
\[
 b_{p,t}=d_p+\sum_{s<t}x_{p,s}.
\]
A prerequisite taken concurrently is insufficient. Missing catalogue leaves
evaluate to false, so an OR expression can use another available alternative.
Normalized trees contain course leaves, constants, AND, OR and $k$-of
operators. A Boolean auxiliary variable represents each compound expression
exactly: for children $v_1,\ldots,v_n$,
\[
 z=1\ \Longleftrightarrow\ \sum_{i=1}^n v_i\geq k .
\]
AND uses $k=n$ and OR uses $k=1$. Both directions of the equivalence are
encoded. If $P_c(t)$ is the resulting prerequisite truth value, enforce
$x_{c,t}\leq P_c(t)$. Unselected courses do not force their prerequisite chains.

For an ordinary co-requisite $p$ of $c$, require
$x_{c,t}\leq b_{p,t}+x_{p,t}$. Ordinary co-requisite lists are empty in this
catalogue; the active CS2101/CS2103T pair
has the stronger new-enrolment rule
\[
 x_{\mathrm{CS2101},t}=x_{\mathrm{CS2103T},t}.
\]
Historical completed records are trusted. The model does not retrospectively
audit a supplied transcript.

\subsection{Preclusions and request constraints}
For a precluded pair $(c,p)$, impose $y_c+y_p\leq1$ when at least one course
is not already completed. A listed preclusion in either direction constrains
the pair. The exception for two completed courses preserves the policy of
trusting historical credit.
Required, forbidden and fixed courses impose, respectively,
\[
 y_c=1\ (c\in R),\qquad
 \sum_t x_{c,t}=0\ (c\in F),\qquad
 x_{c,t}=1\ ((c,t)\in J).
\]
A fixed position refers to the remaining-semester index, not to the normal
semester number. For example, positions in $[2,1,2]$ correspond to normal
Semesters 2, 1 and 2.

These constraints interact: CS4246 requires earlier CS3243 or CS3263, while
mandatory CS2109S precludes CS3243 in the frozen catalogue. CS3263 therefore
provides the usable prerequisite route.

\par
\subsection{Workload and requirement groups}
Semester workload includes every newly enrolled catalogue course:
\[
 \sum_{c\in C}u_c x_{c,t}\leq W,\qquad W=16 .
\]
For a requirement group $g$ with membership $M_g$, define
\[
 U_g=\sum_{c\in M_g}u_c y_c,\qquad N_g=\sum_{c\in M_g}y_c .
\]
Each specified unit or course-count bound constrains $U_g$ or $N_g$.
Unspecified bounds add no constraint. A course counts once within a group;
membership in several groups can satisfy several conditions without creating
extra earned credits.

\begin{center}
\begin{minipage}{\linewidth}
\centering
\captionof{table}{Requirement groups in the frozen model.}
\label{tab:groups}
\begin{tabularx}{\linewidth}{@{}lY@{}}
\toprule
Group & Membership and bound\\
\midrule
Breadth & AI breadth courses and industry credit; at least 32 units.\\
Upper breadth & Breadth members at level 4000 or above; at least 12 units.\\
AI primary & Available AI primaries, including foundation CS2109S; at least three courses.\\
AI upper & AI primaries at level 4000 or above; at least one course.\\
Industry & CP3200 in the selected pathway; 6--12 units.\\
Other CP & Non-industry CP breadth courses; at most 12 units. This set is empty in the catalogue.\\
Programme lower & Programme members at levels 1000--4000; at least 64 units, representing 80\% of the fixed 80-unit target.\\
\bottomrule
\end{tabularx}
\end{minipage}
\end{center}
Foundation and mathematics courses are individually mandatory; their credits
do not also earn breadth credits. CS2109S can satisfy an AI focus count
without contributing breadth units. The 80-unit target is a minimum, not an
exact selected total: six CP3200 units plus four-unit breadth courses require
at least 34 breadth units, yielding at least 82 Programme units. The 64-unit
bound uses the fixed requirement, and permitted overshoot is accepted.

\subsection{Solver implementation and independent checking}
The implementation uses Python 3.11.9 and OR-Tools 9.15.6755. The Boolean
CSP is compiled into CP-SAT with no objective. Unit coefficients are multiplied
by 100 to obtain integers; values that cannot be represented in nonnegative
hundredths are rejected. CP-SAT supports feasibility solving and distinguishes
proven infeasibility from an unknown result \cite{cpsat}. We use one search
worker, random seed 3263 and a ten-second evaluation solver limit.
Backend feasible or optimal statuses are returned as feasible assignments;
an unknown status is never relabelled infeasible.

An independently implemented arithmetic and set-based checker assesses every
feasible output against the \emph{original} specification. It checks schema,
semester count, duplicates, eligibility, offerings, prerequisite order,
co-requisites, preclusions, workload, request constraints and group totals.
It does not invoke the solver to validate a schedule.

The freeze records 25 passing tests covering real-course examples, invalid
schedules, completed credit, logical prerequisites, preclusions, fractional
credits, group accounting and experiment-input integrity.

\par
\section{Benchmark construction}
\subsection{Study stages and sampling}
The benchmark contains 24 development cases and 120 evaluation cases. Four
horizons, $H\in\{7,6,5,4\}$, are crossed with three designed difficulty levels.
Each of the 12 cells contains one feasible and one infeasible development
case, plus seven feasible and three infeasible evaluation cases. Evaluation
therefore includes 84 feasible and 36 infeasible requests.

\begin{center}
\begin{minipage}{\linewidth}
\centering
\captionof{table}{Study-stage profiles used to construct the benchmark.}
\label{tab:stages}
\begin{tabular}{@{}rrrr@{}}
\toprule
\shortstack{Remaining\\semesters} & \shortstack{Past\\semesters} &
\shortstack{First normal\\semester} & \shortstack{Academic courses\\passed}\\
\midrule
7 & 1 & 2 & 2--4\\
6 & 2 & 1 & 5--7\\
5 & 3 & 2 & 8--9\\
4 & 4 & 1 & 10--12\\
\bottomrule
\end{tabular}
\end{minipage}
\end{center}

Synthetic academic histories satisfy earlier prerequisites, offerings,
co-requisites and past workload limits. CP3200 credit is added separately.
The development and evaluation splits have disjoint completed-course sets.
Eight distinct histories are used in development and 70 in evaluation.
Histories recur within each split, and both splits use related request
templates and the same catalogue. These are correlated controlled cases,
not independent student samples or a template-disjoint generalization test.

\begin{center}
\begin{minipage}{\linewidth}
\centering
\captionof{table}{Designed restrictions in feasible requests.}
\label{tab:difficulty}
\begin{tabular}{@{}lrrr@{}}
\toprule
Difficulty & Requested electives & Fixed placements & Forbidden alternatives\\
\midrule
Easy & 1 & 0 & 0\\
Medium & 3 & 2 & 2\\
Hard & 6 & 5 & 4\\
\bottomrule
\end{tabular}
\end{minipage}
\end{center}

The catalogue, programme bounds and workload cap remain fixed. References
are any satisfying assignments; extra permitted courses and empty semesters
are valid. Fixed placements are selected from a feasible reference and
forbidden alternatives from unused courses, followed by revalidation.
The labels describe construction choices, not measured proof difficulty.
More fixed placements may reduce assignment freedom.

\subsection{Feasible references and negative cases}
Every feasible reference passes the independent checker. Infeasible cases
are derived from a checked feasible precursor and carry a concrete
contradiction. The nine annotated mechanisms are: required-but-forbidden
courses; unavailable offerings; prerequisites impossible in the first
remaining semester; preclusions; split simultaneous enrolments; prerequisite
ordering; co-requisite workload overload; insufficient lower-level credits;
and exhaustion of prerequisite alternatives. Each mechanism occurs in four
evaluation cases. Solver-unknown requests are not accepted as infeasible
labels.

A witness is a sufficient contradiction rather than a claim about the sole
cause. For example, an overload request can also place an advanced course
before its prerequisites. This limits conclusions about reasoning difficulty
by mechanism. A valid infeasibility explanation can identify another sufficient
contradiction.

Evaluation cases were selected without language-model answers. A 12-case
paired development pilot achieved 12/12 for both approaches and informed
inspection and scoring; its scores are excluded from evaluation totals.
Source and assumption review was confirmed on 11 October 2026.
Reference-label validation uses the solver, independent
schedule checks and recorded contradiction witnesses.

\par
\section{Evaluation methodology}
\subsection{Paired approaches and information supplied}
Both approaches receive the same natural-language request and frozen domain
facts: normalized course constraints and an enumerated narrative of group
membership and bounds. Gold specifications, reference plans and labels are
excluded. Translation extracts a fixed contract from supplied facts rather
than interpreting unprocessed university policy pages.

In \emph{direct planning}, the model returns a plan JSON object with status,
one course list per remaining semester when feasible, and a brief explanation.
It may return infeasible or unknown. In \emph{solver-assisted planning}, the
model returns a specification containing completed courses, alternating
normal-semester numbers, workload cap, required and forbidden courses, fixed
placements, and requirement groups. A fixed compiler independently enforces
catalogue course constraints. The model generates no executable solver code.
The evaluator then runs CP-SAT. The inference paths are
\[
\begin{aligned}
\text{Direct:}\quad & \text{request + facts}\ \to\ \text{LLM}\ \to\ \text{plan},\\
\text{Assisted:}\quad & \text{request + facts}\ \to\ \text{LLM}\ \to\
\text{specification}\ \to\ \text{CP-SAT}\ \to\ \text{plan}.
\end{aligned}
\]
Offline evaluation checks both plans against the same original specification;
neither model receives checker feedback.

The task-specific instructions differ. In particular, the translation prompt
explicitly states that completed mandatory courses remain required; the
direct prompt prohibits re-enrolment without equally explicit guidance on
the meaning of requested inclusion for a passed course. This distinction is
a potential confound in attributing any observed advantage to the solver.

\subsection{Execution controls and provenance}
The recorded model is \code{gpt-6.1-sol} at medium reasoning effort, run
through subscription-authenticated Codex CLI \code{0.162.0-alpha.17.2}.
Each attempt uses a fresh ephemeral session, schema-constrained output
and a 180-second timeout. Model-side tools are disabled and event records
audited. Cases run in fixed order, direct then assisted, without repairs,
retries or feedback from reference outcomes.

Run \code{csp16-evaluation-v1} freezes settings, case order, scoring and a
426-file snapshot. Completion checks verified all 240 prompt hashes,
frozen files, documentation and pilot evidence. Commit \code{d251777}
preserves the implementation, results and raw responses \cite{repo}.

\subsection{Scoring and secondary analysis}
The primary policy, \code{outcome-v1-explanation-secondary}, accepts a
schema-valid feasible response only when its plan passes the original-specification
checker and the reference case is feasible. An infeasible prediction is
correct when it agrees with a validated infeasible label. Unknown outputs,
execution failures and malformed responses count as unsuccessful completed
attempts. Accuracy uses completed selected attempts.

Explanations do not change a valid schedule's primary score. They are
reviewed separately. Translation auditing first compares unordered field
values and named groups with the reference specification, then reviews
differences as equivalent, non-equivalent or unresolved. A matching solver
outcome alone is insufficient evidence of faithful translation.

Runtime is recorded process time plus solver time for the assisted approach;
shared post-hoc checking and manual review are excluded. CLI token counts
describe usage, not dollar costs. The subsequent diagnostic explanation audit
checks explicit numeric claim patterns.

\par
\section{Results}
\subsection{Primary correctness and execution}
All 240 attempts completed with execution status \code{ok}, and no tool
invocations were recorded. Both approaches correctly identified all 36
infeasible cases. Every returned feasible schedule passed the original
specification checker. The direct approach instead declared one feasible
case infeasible, producing the only primary failure.

\begin{center}
\begin{minipage}{\linewidth}
\centering
\captionof{table}{Primary outcomes on the 120 evaluation cases.}
\label{tab:outcomes}
\begin{tabular}{@{}lrr@{}}
\toprule
Measure & Direct & Solver-assisted\\
\midrule
Valid plans on feasible cases & 83/84 & 84/84\\
Correct infeasibility judgments & 36/36 & 36/36\\
False infeasibility judgments & 1 & 0\\
Correct outcomes overall & 119/120 & 120/120\\
Overall accuracy & 99.17\% & 100.00\%\\
Both-correct paired cases & \multicolumn{2}{c}{119/120}\\
Assisted-only correct paired cases & \multicolumn{2}{c}{1/120}\\
\bottomrule
\end{tabular}
\end{minipage}
\end{center}

Both methods scored 10/10 in eleven of the twelve horizon--difficulty
cells. In the four-semester easy cell, direct planning scored 9/10 and
assistance 10/10. Thus accuracy did not decline with designed difficulty or
longer horizons. The sole discordant pair, analysed in
Section~\ref{sec:false-infeasibility}, gives an observed difference of
0.83 percentage points.

\subsection{Runtime and reported usage}
Mean recorded time was 22.41 seconds for direct planning and 31.60 seconds
for solver assistance: 9.19 seconds, or approximately 41.0\%, longer.
Mean local assisted solver time, including compilation, was approximately
0.105 seconds, about 0.33\% of assisted time.

\begin{center}
\begin{minipage}{\linewidth}
\centering
\captionof{table}{Mean recorded time by reference label, including the
assisted solver stage. All selected attempts are included.}
\label{tab:runtime}
\begin{tabular}{@{}lrrr@{}}
\toprule
Reference label & Cases & Direct (s) & Assisted (s)\\
\midrule
Feasible & 84 & 27.28 & 31.68\\
Infeasible & 36 & 11.04 & 31.42\\
All & 120 & 22.41 & 31.60\\
\bottomrule
\end{tabular}
\end{minipage}
\end{center}

The overhead depends strongly on the reference label: 4.39 seconds
(16.1\%) on feasible cases, versus 20.38 seconds (184.6\%) on infeasible
cases. Direct responses to infeasible cases average 165 output tokens,
compared with 676 on feasible cases; assisted specifications average
1,117 and 1,108 respectively. Short direct contradiction responses and
complete translated specifications are consistent with this timing pattern.

Reported input-token totals were 2,332,737 and 2,359,601; output-token totals
were 62,736 and 133,277 for direct and assisted execution respectively.
These records describe process latency rather than a causal decomposition:
execution order, network and service conditions were not randomized.

\par
\section{Error analysis and translation fidelity}
\subsection{The false infeasibility case}
\label{sec:false-infeasibility}
The request illustrated in Table~\ref{tab:worked},
\code{test-h4-easy-01}, produces the only primary failure. The direct response treats CS1101S
as a compulsory new enrolment and declares a contradiction with the
no-retake rule. The formal interpretation distinguishes
$y_{\mathrm{CS1101S}}=1$ from $x_{\mathrm{CS1101S},t}=1$: past completion
satisfies the requirement, and there is no fixed new placement for this
course. The saved assisted schedule passes all original constraints.
This is an interpretation error under the frozen semantics. The prompt
difference discussed above prevents attributing it solely to a lack of
symbolic search.

\subsection{Canonical differences and equivalent specifications}
Automatic field comparison found 115/120 canonical translations. All five
remaining specifications were reviewed as equivalent under the frozen
catalogue and compiler semantics. Table~\ref{tab:translation} summarizes
the case-specific reasons; each concerns constraints in both directions.
These classifications are recorded separately from the unchanged automated
diagnostic fields \cite{audit}.

\begin{center}
\begin{minipage}{\linewidth}
\centering
\captionof{table}{Semantic review of the five noncanonical translations.
Case identifiers omit the common \code{test-} prefix.}
\label{tab:translation}
\begin{tabularx}{\linewidth}{@{}lY@{}}
\toprule
Case(s) & Constraint-level equivalence evidence\\
\midrule
\code{h6-easy-07} & A new group requires all 14 courses already individually required, repeating the same conditions.\\
\shortstack[l]{\code{h4-easy-03}\\\code{h4-easy-07}} & Four added forbidden courses are already non-selectable; CP3200 past credit remains counted.\\
\code{h4-hard-01} & The 19 individual requirements are replaced by a group containing exactly those 19 members and requiring all of them.\\
\code{h4-hard-08} & Added required CS4278 is already forced by its original fixed placement in position 1.\\
\bottomrule
\end{tabularx}
\end{minipage}
\end{center}

The canonical match rate is 95.83\%, although all 120 specifications are
either canonical matches or reviewed as semantically equivalent. The five
judgments follow from redundant or equivalent constraints, independently
of whether solver outcomes match.

\subsection{Numeric explanations of valid direct plans}
A diagnostic audit checked 453 explicitly matched numerical claims across
all 83 feasible direct responses. Literal-claim review and independent
checker evidence confirmed 17 incorrect claims in 15 responses (18.1\%).
Thirteen overstate AI focus-course counts, including ten AI-upper counts;
the other four concern credit totals or a semester load.
The affected plans still satisfy all original constraints.

\begin{center}
\begin{minipage}{\linewidth}
\centering
\captionof{table}{Examples of explanation errors in independently valid plans.}
\label{tab:explanations}
\begin{tabularx}{\linewidth}{@{}lYll@{}}
\toprule
Case & Claim type & Reported & Checked\\
\midrule
\code{h4-easy-04} & AI primary / AI upper counts & 8 / 4 & 7 / 3\\
\code{h4-hard-02} & Final-semester workload & 16 units & 12 units\\
\code{h6-hard-03} & Programme-lower credits & 70 units & 74 units\\
\bottomrule
\end{tabularx}
\end{minipage}
\end{center}

The audit is not exhaustive: only defined numeric patterns were checked, so
18.1\% is the observed proportion with detected errors, not an estimate of all
possible prose inaccuracies. No primary score was changed. Assisted outputs
use a generic solver status explanation; this study does not compare two
equally detailed natural-language explanation systems.

\par
\section{Discussion and threats to validity}
\subsection{What the comparison establishes}
With the supplied catalogue and structured facts, direct planning was
already highly accurate. Assistance added one correct outcome and an
inspectable specification, at greater process latency and output-token
usage. That latency difference was concentrated in infeasible requests;
local solving itself was inexpensive. The experiment therefore distinguishes
the cost of obtaining a formal specification from the cost of solving it.

The clearest additional finding is inaccurate numerical explanation
despite valid planning. Most observed errors overcount focus membership,
so an assistant should compute displayed counts and credits from the
checked schedule and completed record. This is a design implication;
an improved explanation pipeline has not been evaluated here.

The translation audit shows why literal matching is insufficient:
five different representations express equivalent constraints. Conversely,
a solver validates the specification it receives, so checking schedules
against the original request remains necessary. Faithful interpretation,
valid assignments and accurate explanations are separate engineering obligations.

\subsection{Representation size and designed difficulty}
With 87 courses and $H$ remaining semesters, the main model contains $87H$
enrolment variables and 87 completion/selection variables before logical
auxiliaries. Ignoring other constraints, at-most-once enrolment permits at
most $(H+1)^{87}$ course-position combinations; completed credit, eligibility
and offerings substantially reduce this space. Saved validation examples
contain 2,160 backend variables and 3,516 constraints for seven semesters,
and 1,248 variables and 2,127 constraints for four semesters. These are
representation measurements of two examples, not counts of explored states
or evaluation-case averages.

Large formal representation size did not translate into observed answer
failures on the designed hard cells. Many negative cases contain a short
sufficient contradiction, while fixed placements constrain choices.
The benchmark therefore probes specified constraint interactions without
locating a general difficulty threshold at which direct planning fails.
No search-state enumeration, optimization comparison or component ablation
was performed.

\subsection{Limits of the evidence}
Cases are conditioned on solver validation and some restrictions derive
from reference plans. The 70 evaluation histories recur in related templates,
so the cases are not independent student samples. Disjoint completed-course
sets separate the splits, but shared construction limits generalization.
Supplied group memberships and normalized facts also make interpretation
easier than reading unprocessed university policies.

Only one model setting and one attempt per case were tested. The single
discordant pair, prompt asymmetry and fixed execution order limit claims
about general solver advantages, repeatability and causal latency.
The solver and checker share source interpretations. The Codex-supported
semantic and numeric audits preserve case-level evidence without constituting
an independent human replication.

\par
\section{Responsible AI considerations}
False infeasibility can reject a valid study route, while inaccurate
explanations can create misleading confidence in a valid schedule.
A deployment should show checked constraint evidence, distinguish unknown
from infeasible, and route disputed interpretations to university guidance.

The benchmark uses synthetic histories and public course information.
Deployment with personal transcripts would require consent, access controls,
data minimization and a policy for missing or disputed completed credit.

Standard-entry assumptions and internship credit exclude some pathways;
performance for bridging requirements, exemptions, exchange, double degrees
or different internship opportunities is untested. These inputs require
explicit handling. The workload cap and repeated offerings are modelling
choices, and feasibility does not certify graduation eligibility or future
course availability.

Retained source snapshots, interpretations, responses, hashes and independent
checks support traceability. Displayed totals should come from checked
evidence, and interpreted constraints should remain inspectable. AI assistance
is disclosed because Codex supported both implementation and secondary review.

\section{Conclusion}
The evaluation found 119/120 correct direct outcomes and 120/120 assisted
outcomes, with approximately 41\% greater mean recorded time for assistance.
Both correctly identified all 36 infeasible cases; every returned feasible
schedule passed the original checker. Five noncanonical translations were
equivalent, while 15 valid direct responses contained numeric explanation errors.

For these structured instances, direct planning combines strong accuracy
with faster responses, especially for infeasible requests. Assistance
provides inspectable constraints and correct outcomes at higher process
cost, with local solving contributing little to that cost. Numerical
explanations remain a distinct reliability concern. Evaluating schedule
validity, translation fidelity and explanation accuracy separately therefore
gives a more useful account of an assistant than primary accuracy alone.

\section{Team contributions and AI assistance}
The following table summarizes the team's division of work.
\begin{center}
\begin{minipage}{\linewidth}
\centering
\begin{tabularx}{\linewidth}{@{}lY@{}}
\toprule
Team member & Contribution\\
\midrule
Brandon Chong Wee Kiat & Formulated the project idea; collected and processed the initial dataset.\\
Lu Ruijian & Built the benchmark; implemented the solver-assisted pipeline; coordinated and organized the report.\\
Jonathan Graziano Yang & Conducted the evaluation; contributed to report writing.\\
\bottomrule
\end{tabularx}
\end{minipage}
\end{center}
Codex supported source acquisition and normalization, CSP and checker
implementation, benchmark generation, execution tooling, test construction,
experiment execution, secondary analysis and report preparation, including
revision informed by three separate AI reviewers. Student participants
selected the scope and modelling policies, directed revisions and confirmed
completion of the source/assumption review. The team is responsible for
the technical interpretation and conclusions.

\clearpage
\typeout{REPORT_BODY_PAGES=\number\numexpr\value{page}-1\relax}
\pagenumbering{roman}
\begin{thebibliography}{9}
\raggedright

\bibitem{llmp}
B. Liu, Y. Jiang, X. Zhang, Q. Liu, S. Zhang, J. Biswas and P. Stone.
\emph{LLM+P: Empowering Large Language Models with Optimal Planning Proficiency}.
arXiv:2304.11477, 2023.
\url{https://arxiv.org/abs/2304.11477}.

\bibitem{logiclm}
L. Pan, A. Albalak, X. Wang and W. Y. Wang.
\emph{Logic-LM: Empowering Large Language Models with Symbolic Solvers
for Faithful Logical Reasoning}.
arXiv:2305.12295, 2023.
\url{https://arxiv.org/abs/2305.12295}.

\bibitem{nuscohort}
NUS School of Computing.
\emph{CS Cohort 2026/2027: Degree Requirements}.
\url{https://www.comp.nus.edu.sg/cug/per-cohort/cs/cs-26-27/}.
Experiment source snapshot collected 9 October 2026.

\bibitem{nusfocus}
NUS School of Computing.
\emph{CS Focus Areas: Artificial Intelligence}.
\url{https://www.comp.nus.edu.sg/programmes/ug/focus/}.
Experiment source snapshot collected 9 October 2026.

\bibitem{nusmods}
NUSMods.
\emph{Public API, academic year 2026--2027}.
\url{https://api.nusmods.com/v2/}.
Per-course endpoint:
\code{/2026-2027/modules/COURSE.json}.
Experiment source snapshots collected 9 October 2026.
Project and API provenance:
\url{https://github.com/nusmodifications/nusmods}.

\bibitem{cpsat}
Google OR-Tools.
\emph{CP-SAT Solver}.
\url{https://developers.google.com/optimization/cp/cp_solver}.
Accessed 14 October 2026. Experiment software version: 9.15.6755.

\bibitem{repo}
CS3263 project team.
\emph{AI Degree Planning Assistant: frozen implementation and paired evaluation}.
Repository commit \code{d2517771fd28edf657475a4ea1f68096458e035b}.
\url{https://github.com/RuijianLu/CS3263Project/tree/d2517771fd28edf657475a4ea1f68096458e035b}.
Primary evidence: \code{results/csp16-evaluation-v1/metrics.json},
\code{summary.csv} and \code{completion.json};
freeze and raw records: \code{results/runs/csp16-evaluation-v1/}.

\bibitem{audit}
CS3263 project team, with Codex assistance.
\emph{Secondary audit of csp16-evaluation-v1}.
\url{https://github.com/RuijianLu/CS3263Project/blob/d2517771fd28edf657475a4ea1f68096458e035b/results/csp16-evaluation-v1/manual_audit.json}.
Numeric-claim evidence is preserved in
\code{results/csp16-evaluation-v1/explanation\_numeric\_checks.json}
at the same commit.

\end{thebibliography}
\end{document}
